\section{The Front End}

\subsection{Structure}

The interface is a single page built with Svelte 5 (runes), TypeScript and Vite.  Its
logic is kept in small \emph{pure} modules, which import nothing from the browser, so that
they run under Node's own test runner: ordering (\code{sort}), filtering (\code{filter}),
column layout (\code{columns}), preferences (\code{prefs}), preview classification and text
helpers (\code{preview}), the path and location functions (\code{format}, \code{goto}),
the key mapping (\code{keys}), the Markdown conversion (\code{markdown}), and where the page was
loaded from (\code{base}).  The components hold only what needs the DOM.

\subsection{Decisions}

\begin{description}
  \item[Virtualized list.]  Only the rows on screen are drawn, from a scroll position and a
        fixed row height; a folder of 100{,}000 entries draws about 30 rows.
  \item[Location as state.]  The current folder, and the entry to select, are in the URL
        fragment (\code{\#/path?select=name}), which is never sent to the server, so folders
        are bookmarkable and Back works without any routing library.  Parsing is strict: a
        bad escape is not silently replaced.
  \item[Requests relative to the page.]  Every request is formed from the page's own
        directory, so the random prefix (which scopes the cookie) needs no configuration.
  \item[Sorting in the browser.]  The server streams entries in the order the file system
        gave them; the browser sorts, with a comparison that is a total order and does not
        depend on the input order (a test).  A time that cannot be parsed sorts as the oldest,
        so the order stays total.
  \item[Names are displayed, not trusted.]  Characters that would disguise a name are replaced by
        visible markers for display only.  Names are always inserted as text, never as markup.
  \item[Preferences are optional.]  They are kept in local storage with every read guarded, and
        every value falls back on its own.  Local storage is per origin, which includes the
        port; that is why the server remembers its port (Section~\ref{sec:port}).
  \item[The pane grows with the window.]  \code{clampPaneWidth} bounds the preview pane
        between 240 pixels and the window's width less \code{RESERVED\_FOR\_LIST}
        (320 pixels, matched by a CSS \code{max-width}), and the page tracks the window's
        width live, so the pane can widen with the window and is narrowed if the window
        shrinks.  An earlier fixed ceiling of 900 pixels was reported by the user as a
        usability limit, not found by a test.
  \item[Every letter types ahead.]  Search, Go to path and Copy are on Alt+S, Alt+G and
        Alt+C, read by physical key (\code{KeyboardEvent.code}) because macOS changes the
        character Option produces; so no plain letter is reserved and type-ahead
        (\fq{KEY-3}) can reach every name.
  \item[A click selects.]  A plain left click on a file name calls \code{preventDefault}, so
        it selects the row instead of following the link's \code{download} attribute;
        modified clicks and the context menu still reach the browser's own download
        (\fq{MOUSE-1}).
  \item[No dynamic code.]  There is no use of \code{eval}, of \code{innerHTML}, or of any
        other markup sink except one: the output of the syntax highlighter, which escapes its
        input.
\end{description}

\subsection{The preview pane as a state machine}

For a selected entry the pane is in one of a few states (idle, loading, folder, image, Quick
Look picture, PDF, archive, text, binary, empty, cloud-only, error).  A file name gives a \emph{hint}, the
server decides from the bytes, and a hint the server rejects falls back to the ordinary head
request, which never falls back again, so every run ends in a terminal state after at most two
requests.  (A Quick Look picture that fails falls back to that head request, or, for a package
document, which lists as a folder, straight to the folder state.)  Every request started for a selection, including a fall back, shares one
cancellation handle and is cancelled when the selection changes, so a late answer can never be
shown for the wrong file.  The laws are in the algebraic specification.

The Details panel (metadata and extended attributes) is a second, independent request
against the same selection, and keeps its own error state rather than sharing the main
preview's.  An earlier version used one shared variable for both: whichever request
settled last silently overwrote the other's message, and the panel's own error was
suppressed whenever the main preview had already failed, which could leave it reading
``Loading…'' forever even though its own request had already come back (Finding~18).

The folder listing has its own, separate error mapping (\code{describe} in
\code{App.svelte}), independent of the preview pane's, and it had never been given a case
for status 409: a cloud-only folder (Finding~19) fell through to the generic ``Server
error (\emph{n}).'' instead of the same cloud message the preview pane already shows for
a cloud-only \emph{file} (Finding~20).  Two error mappings for the same status code, one
of them incomplete, is the same shape of gap as Findings 4 through 6: a case a status code
can take that nothing had exercised.
Two independent outcomes need two independent variables.

\subsection{Hover peek}

A short delay before the request, a cache, a limit on simultaneous requests, and cancellation
on pointer movement, scroll or key press.  It requests the same head endpoint as the preview,
so it is subject to the same rules, and is off if the user turns it off.
